\setcounter{section}{1}\setcounter{subsection}{0}
% Original source lines 90--204
\subsection{Notations and conventions}
In this article, $M$ denotes a $d$-dimensional (with $d \geq 3$) smooth manifold endowed with a metric tensor $g$ or conformal class of metrics $\cc$, both of which will be positive definite. Furthermore, $\Lambda$~denotes a $(d-k)$-dimensional smooth submanifold of $M$, and $0 < k < d$. Throughout the entirety of this article, we work only locally and so will assume that trivializations are always available for any section of any bundle. In the Riemannian setting $(M,g)$, we denote the Levi-Civita connection by $\nabla$. Our convention for the curvature tensor is given~by
\[
R(x,y)z = (\nabla_x \nabla_y - \nabla_y \nabla_x)z - \nabla_{[x,y]} z .
\]
Here, $x$, $y$, $z$ are vector fields on $M$ and $[x,y]$ is the Lie bracket of $x$ and $y$. In general, for any bundle $\mathcal{V}M$ over $M$, we will denote sections of this bundle by $\Gamma(\mathcal{V} M)$. That is, $x,y,z \in \Gamma(TM)$. When necessary, we will refer to arbitrary tensor products of a bundle (and its dual) using capital Greek letter, so that, for example, $T^\Phi M$ is some unspecified tensor product of $TM$ and~$T^* M$.

For ease of calculation, we will use Penrose's abstract index notation, labelling such indices with letters from the beginning of the Latin alphabet $a$, $b$, $c$, etc., so that, for example, the~Riemann curvature tensor may be expressed as
\[
R_{ab}{}^c{}_d x^d = [\nabla_a, \nabla_b] x^c .
\]
Using the metric and its inverse, indices can be raised and lowered, as usual; the Einstein summation convention is used to indicate contraction. (Sometimes, when the notation is unambiguous, we will use dot product notation to represent contraction; for example, when $\omega$ is some 1-form, we will sometimes write $\nabla^a \omega_a =: \nabla \cdot \omega$.) In this language, we have $\mathrm{Ric}_{ab} := R_{ca}{}^c{}_b$ and its corresponding scalar is $\mathrm{Sc} := \mathrm{Ric}_a^a$. When $d > 2$, a particular trace-correction of the Ricci tensor, known as the Schouten tensor, arises frequently in conformal considerations,
\[
P_{ab} := \frac{1}{d-2} \left(\mathrm{Ric}_{ab} - \frac{1}{2(d-1)} \mathrm{Sc}\, g_{ab} \right) .
\]
The trace of this tensor is denoted by $J := P_a^a$. Written in terms of the Riemann tensor and the Schouten tensor, we may express the Weyl tensor as
\[
W_{abcd} = R_{abcd} - g_{ac} P_{bd} + g_{bc} P_{ad} + g_{ad} P_{bc} - g_{bd} P_{ac} .
\]
In dimensions $d > 2$, the Cotton tensor is given~by the covariant curl of the Schouten tensor,
\[
C_{abc} = \nabla_a P_{bc} - \nabla_b P_{ac} .
\]
When $d > 3$, this can be reexpressed in terms of the divergence of the Weyl tensor,
\[
C_{abc} = \frac{1}{d-3} \nabla^d W_{dcab} .
\]

In the context of submanifold embeddings, we will use Greek indices to represent tensors on normal frame bundles. Objects that are along the submanifold will be denoted by overbars~$\bar{\bullet}$, so that, for example, $\bar{\nabla}$ represents the induced Levi-Civita connection on $(\Lambda,\bar{g})$, where $\bar{g}$ is the induced metric from $\Lambda \hookrightarrow (M,g)$. In general for any index notation, we will use round brackets~$(\cdots)$ to indicate unit-normalized symmetrization and square brackets $[\cdots]$ to indicate unit-normalized antisymmetrization. Furthermore, we will use a subscript $\circ$ to indicate the tracefree part of a particular symmetric index configuration. However, because Latin and Greek indices represent different types of objects, we may sometimes write $(\alpha a b \beta)$ to specifically mean the symmetrization over $\alpha$ and $\beta$, and similarly for antisymmetrization. To be explicit, whichever type of index is closest to the opening and closing brackets is the type of index that is (anti)symmetrized over. On the other hand, when excluding an index from (anti)symmetrization is necessary, we will use vertical bars, so that $(\alpha | \beta | \gamma)$ refers to the symmetrization over~$\alpha$ and~$\gamma$. Note that Greek indices will always be subscripts, so repeated Greek indices will denote summation in the usual way.

\section[Orthonormal frames for normal bundles of Riemannian submanifold embeddings]{Orthonormal frames for normal bundles \\ of Riemannian submanifold embeddings} \label{riem-frames}

We begin with the case of smooth Riemannian submanifold embeddings $\iota \colon \Lambda^{d-k} \hookrightarrow (M,g)$ with $d \geq 3$. For notational convenience, we will often view $\Lambda$ as a subset of $M$. We now review classical submanifold geometry.

\subsection{Classical submanifold geometry}

 At a point $p \in \Lambda$, every tangent space of the bulk space along $\Lambda$ can be decomposed according~to
\begin{align} \label{TM-decomp}
T_p M = T_p \Lambda \oplus N_p \Lambda ,
\end{align}
where $N_p \Lambda$ is the normal space at $p \in \Lambda$. The normal bundle over $\Lambda$ is defined as $N \Lambda := \sqcup_{p \in \Lambda} N_p \Lambda$.

As the pullback map is well defined, we can construct the \textit{intrinsic metric}, also sometimes called the \textit{first fundamental form}, via the pullback: $\bar{g} := \iota^* g$. Note that this metric is compatible with the decomposition above, and given an orthornormal frame $\{n_{\alpha}\}$ for $N \Lambda$, we may express this invariant in index notation according to $\bar{g}_{ab} = g_{ab}|_{\Lambda} - n_{a \alpha} n_{b \alpha}$. In the notation $n^a_{\alpha}$ (and $n_{a \alpha}$), the Latin index is an abstract index whereas the Greek letter indexes some particular element of the basis -- this will always be the case. Closely related to the first fundamental form is the projection operator given~by $\bar{g}^a_b := \delta^a_b - n^a_{\alpha} n_{b \alpha}$.

With a metric in hand, one may construct the Levi-Civita connection on $\Lambda$, denoted $\bar{\nabla}$. A~classical result of submanifold geometry is that this connection is directly obtainable from the Levi-Civita connection $M$ and the projection operator,
\[
\bar{\nabla}_{\bar v} \bar{u}^a = \bar{g}^a_b \nabla_{\iota_* \bar{v}} (\iota_* \bar{u})^b ,
\]
where $\bar{v},\bar{u} \in \Gamma(T \Lambda)$ and $\iota_* \bar{v}$, $\iota_* \bar{u}$ are their respective pushforwards. Note that we will often write~$\nabla^\top$ for the pullback of the bulk Levi-Civita connection -- this way, we may use abstract index notation and we do not need to carry around arbitrary vector fields. In a similar way, we use $\top$ both as an operator and as a superscript to denote the projection to the submanifold, e.g., we will write $P^\top$ to represent the projection of the Schouten tensor of $(M,g)$ to $\Lambda$.

If we instead project $\nabla_{\iota_* \bar{v}} (\iota_* \bar{u})$ along a unit normal vector, we obtain a classical result: the~\textit{second fundamental form}. Indeed, relative to an orthonormal frame on $N \Lambda$, we define
\[
\bar{u}^a \bar{v}^b \II_{ab \alpha} := -n_{a \alpha} \nabla_{\iota_* \bar{v}} (\iota_* \bar{u})^a .
\]
A short calculation shows that this can be expressed in a simpler way according to
\[
\II_{ab \alpha} = \bar{g}_{bc} \nabla^\top_a n^c_{\alpha} \in \Gamma\bigl(\odot^2 T^* \Lambda\bigr) .
\]
It is worth pointing out that in our calculations, we view $\II_{ab \alpha}$ as a tuple of 2-tensors depending on an orthonormal frame, as per the above section space. However when an orthonormal frame is not prescribed, the second fundamental form is still well defined, taking values in the normal bundle $N \Lambda$.
Accordingly, one may express the Levi-Civita connection of $M$ restricted to $\Lambda$ via the following identity:
\[
\nabla_{\iota_* \bar{v}} (\iota_* \bar{u}) = \bar{\nabla}_{\bar{v}} \bar{u} - \II_{\alpha}(\bar{u},\bar{v}) n_{\alpha} .
\]

The first and second fundamental forms are sufficient to completely describe an embedded hypersurface in $\mathbb{E}^d$ up to Euclidean motions. This is because, for $k=1$, the normal bundle is one-dimensional and hence (up to orientation) each normal space has a canonical basis given~by $n \in N_p \Lambda$, where $g(n,n) = 1$. (Note that as we work locally for the entirety of this article, we will neglect any issues that arise from orientations.) This rigidity effectively removes all geometric structure from the normal bundle. Unfortunately, this rigidity is lost entirely when we consider higher-codimension submanifolds. Indeed, while orthonormality can certainly be retained, there is (a priori) no uniquely preferred basis for $N \Lambda$, even locally. In what follows, we will find that there is sometimes a preferred basis up to constant ${\rm SO}(k)$ rotations -- but in general, it is not that clear that any preferred basis can be constructed. We will now provide a construction for the final piece of data required to describe embedded submanifolds up to Euclidean motions and also provide notation and language that will be necessary for later considerations.

Just as one can induce a connection on the tangent bundle to $\Lambda$, one may also induce on the normal bundle a connection
\[
D \colon \ \Gamma(N \Lambda) \rightarrow \Gamma(T^* \Lambda \otimes N \Lambda)
\]
given~by
\[
D_{\bar v} \xi = {\perp} (\nabla_{\iota_* \bar{v}} \xi) ,
\]
where $\perp$ is the orthogonal projection to $N \Lambda$, $\bar{v} \in \Gamma(T \Lambda)$ is any vector field along $\Lambda$ and $\xi \in \Gamma(N \Lambda)$ is any normal vector field. Note that this connection is compatible with the metric on the normal bundle given~by $\eta = g|_{\Lambda} - \iota^* g$. Then, for any frame $\{n_{\alpha}\}$ for $N \Lambda$, the connection coefficients are defined by
\[
D_a n^b_{\beta} = \omega_{\gamma a \beta} n^b_{\gamma} .
\]
Given any orthonormal frame $\{n_{\alpha}\}_{\alpha = 1}^k$ for the normal bundle $N \Lambda$, $\omega$ is given precisely by the so-called \textit{normal fundamental forms}~\cite{Spivak} (sometimes called the H{\'a}j{\'i}{\v{c}}ek one-form in the physics literature~\cite{hajicek1975, hajicek1973}):
\[
\omega_{\alpha a \beta} = \beta_{a \alpha \beta} := n^b_{\alpha} \nabla^\top_a n_{b \beta} .
\]
Note that we use $\beta$ only for orthonormal frames.

Associated to the normal bundle is the orthonormal frame bundle $F(N \Lambda)$, which is a principal ${\rm O}(k)$-bundle. So in particular, any two orthonormal frame fields are related to one another by the action of a section $m \in C^\infty(\Lambda, {\rm O}(k))$. Furthermore, the connection coefficients (and thus the normal fundamental forms) transform according to
\begin{align} \label{beta-transf}
\beta \mapsto m \beta m^{-1} + m \mathrm{d} \bigl(m^{-1}\bigr) ,
\end{align}
where adjacency is understood to imply matrix multiplication.
Note that for constant rotations, i.e., when $m$ is a constant matrix, $\beta$ transforms tensorially.

Because $\beta$ is not covariant under the group action, it is not a true invariant of the normal bundle (in that it does not transform linearly under the group action). Rather, its curvature 2-form
\begin{align} \label{norm-curv}
\mathcal{R} = \mathrm{d} \beta + [\beta, \beta]
\end{align}
does transform linearly under the group action and hence \textit{is} a true invariant of the normal bundle. In terms of the connection itself, the curvature 2-form, hereafter known as the \textit{normal curvature}, obeys the standard relation
\[
\mathcal{R}(\bar{u},\bar{v}) \xi = D_{\bar{u}} (D_{\bar{v}} \xi) - D_{\bar{v}} (D_{\bar{u}} \xi) - D_{[\bar{u},\bar{v}]} \xi
\]
for any $\bar{u},\bar{v} \in \Gamma(T \Lambda)$ and $\xi \in \Gamma(N \Lambda)$. The normal curvature is the final ingredient necessary to completely specify an embedded submanifold in $\mathbb{E}^d$ up to Euclidean motions.

To summarize, the three classical invariants of submanifold embeddings are $\bar{g}$, $\II$, and $\mathcal{R}$. It~is these tensors that appear in classical submanifold identities as described by Spivak~\cite{Spivak}, and given below for later use. First is the Gauss equation, relating the tangential component of the Riemann curvature tensor to that of the submanifold $\Lambda$,
\begin{align} \label{gauss-equation}
R_{abcd}^\top = \bar{R}_{abcd} + \II_{ad \alpha} \II_{bc \alpha} - \II_{bd \alpha} \II_{ac \alpha} ,
\end{align}
where $R^\top := \top (R)$, the projection of the Riemann curvature tensor to the submanifold. Similar notation will be used throughout.
Next is the Codazzi--Mainardi equation, relating another component of the Riemann curvature tensor to extrinsic invariants,
\begin{align} \label{codazzi}
R_{abc \alpha}^\top \= 2\bar{\nabla}_{[a} \II_{b]c \alpha} + 2\II_{c[b \beta} \beta_{a] \alpha \beta} .
\end{align}
Note that here, the right-hand side does not look explicitly tensorial -- however, a short calculation shows that it is indeed independent of the choice of orthonormal basis $\{n_{\alpha}\}$. Finally, we have the Ricci equation (which is trivial for hypersurfaces), relating another projection of the Riemann curvature to the normal curvature and the second fundamental form,
\[
R_{ab \alpha \beta}^\top = \mathcal{R}_{ab \alpha \beta} + \II^c_{a \beta} \II_{cb \alpha} - \II^c_{b \beta} \II_{ca \alpha} .
\]



% Original source lines 274--395
While the existence of a ``clean'' gauge choice can sometimes be taken for granted, uniqueness is much harder to demonstrate -- and indeed, it is not always present (see the discussion about the Gribov ambiguity above). Regardless, we can still study the extrinsic geometry of submanifold embeddings equipped with an orthonormal basis for the normal bundle. When a global orthonormal frame exists for $N \Lambda$, it is called a \textit{parallelization} of $N \Lambda$. Since we work locally, we can always assume that an orthonormal frame is specified for $N \Lambda$. We call such an embedding equipped with an orthonormal frame for the normal bundle a \textit{parallelized $($Riemannian$)$ submanifold embedding}, specified by the pair $(\Lambda \hookrightarrow (M,g), \{n_{\alpha}\})$.

\section[Defining maps for parallelized Riemannian submanifold embeddings]{Defining maps for parallelized Riemannian submanifold\\ embeddings} \label{unit-def-funcs}

Given a Riemannian submanifold embedding $\iota \colon \Lambda^{d-k} \hookrightarrow \bigl(M^d,g\bigr)$, we are interested in studying the embedding by examining the jets of a geometrically-adapted \textit{defining map} $s_{\alpha} \colon M \rightarrow \mathbb{R}^k$, which is a $k$-tuple of functions such that for any $p \in \Lambda$, we have that $s_{\alpha}(p) = 0$ for every $\alpha \in \{1, \ldots, k\}$ and ${\rm d}s_{1}(p) \wedge \cdots \wedge {\rm d}s_{k}(p) \neq 0$. That such families exist is well known, at least locally: for every $p \in \Lambda$, there exists a neighborhood $U \subset M$ such that $U \cap \Lambda$ is the zero locus of some defining map $s_{\alpha} \colon U \rightarrow \mathbb{R}^k$. However, there are many such defining maps, and in order for the jets of such a defining map to contain information about the embedding $\iota$, it must be determined canonically by the geometry. To realize this canonical defining map, we consider the \text{Gram matrix} $G_{\alpha \beta} := g^{-1}({\rm d}s_{\alpha},{\rm d}s_{\beta})$ of a given defining map $s_{\alpha}$. In the case of hypersurfaces, we canonically choose $G = \Id$, which is equivalent to the condition that $|{\rm d}s|^2 = 1$. Similarly, in the case of a higher-codimension submanifold we would like to achieve $G= \Id$. Because we will work locally, it is sufficient for our purposes to find a defining map $s_{\alpha}$ such that $G_{\alpha \beta} = \delta_{\alpha \beta} + \mathcal{O}(s^\infty)$, meaning we wish to find a defining map such that for any $m \in \mathbb{Z}_{\geq 1}$, we have that $G_{\alpha \beta} = \delta_{\alpha \beta} + \mathcal{O}(s^m)$. We call such a defining map a \textit{unit defining map}. Going forward, we will abuse notation by also denoting $g^{-1}({\rm d}s_{\alpha},\cdot)$ by $n_{\alpha}$. Note that, when clarification is necessary, we will decorate $n_{\alpha}$ with abstract indices (such as $n^a_{\alpha}$ and $n_{a \alpha}$) to indicate the whether we are referring to the vector or its metric dual.

\subsection{First order} \label{fo-Riem}
Given the parallelization $\{n_{\alpha}\}$ of $N \Lambda$ and any defining map $\hat{s}_{\alpha}$ for $\Lambda$, there exists a linear combination of components of $\hat{s}_{\alpha}$ denoted by $s_{\alpha}$ such that $g^{-1}({\rm d}s_{\alpha},\cdot) \= n_{\alpha}$. This linear combination provides a unique $s_{\alpha}$ modulo terms that are at least second order in a defining function. The behavior of ${\rm d}s_{\alpha}$ away from $\Lambda$, however, is harder to control. Our goal is to obtain a canonical family of unit defining functions adapted to our parallelization satisfying $g^{-1}({\rm d}s_{\alpha},{\rm d}s_{\beta}) = \delta_{\alpha \beta}$. However, this will not be possible in general. Instead, we will choose $s_{\alpha}$ to, in some sense, get as close to the identity Gram matrix condition as possible.
We begin with an iterative procedure. Because the parallelization of $N \Lambda$ is defined to be orthonormal, we can certainly obtain $G_{\alpha \beta} \= \delta_{\alpha \beta}$, so we assume this going forward. This constraint defines a unique family of defining maps that are aligned to the parallelization of $N \Lambda$ and whose elements have a difference given~by \smash{$s'_{\alpha} - s_{\alpha} = \mathcal{O}\bigl(s^2\bigr)$}. For any of these defining maps, it thus follows that
\begin{align} \label{0-order}
G_{\alpha \beta} = \delta_{\alpha \beta} + F^{(1)}_{\alpha \beta \gamma} s_{\gamma} ,
\end{align}
where $F^{(1)}_{\alpha \beta \gamma}$ represents an array of smooth functions on $M$. Our goal will be to pare down this family of defining maps by controlling \smash{$F^{(1)}$}.

However, note that equation~(\ref{0-order}) does not uniquely specify \smash{$F^{(1)}_{\alpha \beta \gamma}$} as an array of functions on $C^\infty M$. To see this, consider the mapping
\[
F^{(1)}_{\alpha \beta \gamma} \mapsto F^{(1)}_{\alpha \beta \gamma} + F^{(1,1)}_{\alpha \beta \gamma \delta} s_{\delta} ,
\]
where \smash{$F^{(1,1)}$} is any array of functions so that \smash{$F^{(1,1)}_{\alpha \beta \gamma \delta} = F^{(1,1)}_{\alpha \beta [\gamma \delta]}$}. This mapping clearly preserves equation~(\ref{0-order}), and so we see that \smash{$F^{(1)}$} is only uniquely determined along $\Lambda$. Nonetheless, we \textit{may} choose \smash{$F^{(1)}$} in a collar neighborhood of $\Lambda$ in a non-canonical way. Going forward, we will assume such a choice has been made; later, we will show that this choice actually yields a unique result.

Because of the symmetry of $G_{\alpha \beta}$, we have that \smash{$F^{(1)}_{\alpha \beta \gamma} = F^{(1)}_{(\alpha \beta) \gamma}$}, and thus as a representation of $S_k$, we can see that the tensor structure of \smash{$F^{(1)}$} in the normal bundle takes the form
\[
\raisebox{3pt}{\scalebox{0.4}{\ydiagram{2,1}}} \; \oplus \; \scalebox{0.4}{\ydiagram{3}} .
\]
For what follows, we will abuse notation and write, for example, that $F^{(1)} \in C^\infty (M,\raisebox{3pt}{\scalebox{0.2}{\ydiagram{2,1}}} \; \oplus \; \raisebox{2pt}{\scalebox{0.2}{\ydiagram{3}}})$.

Now observe that under the mapping
\[
s_{\alpha} \mapsto \tilde{s}_{\alpha} := s_{\alpha} + A^{(1)}_{\alpha \gamma_1 \gamma_2} s_{\gamma_1} s_{\gamma_2} ,
\]
where $A^{(1)}$ is a list of smooth functions on $M$, we have that the zeroth order part of the Gram matrix condition given in equation~(\ref{0-order}) is preserved while changing $F^{(1)}$. Indeed, by inspection we may consider only those \smash{$A^{(1)}$} satisfying \smash{$A^{(1)}_{\alpha \gamma_1 \gamma_2} = A^{(1)}_{\alpha (\gamma_1 \gamma_2)}$}, and so $A^{(1)} \in C^\infty (M , \raisebox{3pt}{\scalebox{0.2}{\ydiagram{2,1}}} \; \oplus \; \raisebox{2pt}{\scalebox{0.2}{\ydiagram{3}}})$. Evidently, it may be possible to choose \smash{$A^{(1)}$} so that it cancels \smash{$F^{(1)}$}.

Differentiating the above display and using the definition of the Gram matrix for the new defining map (combined with equation~\eqref{0-order}), we obtain
\begin{align}
\tilde{G}_{\alpha \beta} &{}= \delta_{\alpha \beta}+ \bigl(F^{(1)}_{\alpha \beta \gamma} + 4 A^{(1)}_{(\alpha \beta) \gamma}\bigr) s_{\gamma} \nonumber
\\
&\quad{}+ \bigl(4 A^{(1)}_{\alpha (\rho \gamma_1)} A^{(1)}_{\beta (\rho \gamma_2)} + 2 n_{a (\alpha} \nabla^a A^{(1)}_{\beta) (\gamma_1 \gamma_2)} + 4 A^{(1)}_{\alpha (\rho \gamma_1)} F^{(1)}_{\beta \rho \gamma_2} \bigr) s_{\gamma_1} s_{\gamma_2}
\nonumber\\
&\quad{}+ \bigl(4A^{(1)}_{\alpha (\rho \gamma_1)} \nabla_{n_{\rho}} A^{(1)}_{\beta \gamma_2 \gamma_3} + 4 A^{(1)}_{\alpha (\omega_1 \gamma_2)} A^{(1)}_{\beta (\omega_2 \gamma_2)} F^{(1)}_{\omega_1 \omega_2 \gamma_3} \bigr) s_{\gamma_1} s_{\gamma_2} s_{\gamma_3} \nonumber\\
&\quad{}+ \bigl(\nabla^a A^{(1)}_{\alpha (\gamma_1 \gamma_2)}\bigr) \bigl(\nabla_a A^{(1)}_{\beta (\gamma_3 \gamma_4)} \bigr) s_{\gamma_1} s_{\gamma_2} s_{\gamma_3} s_{\gamma_4} ,
\label{1-order-correction}
\end{align}
where the above is symmetric in $\alpha \beta$ (even when the round brackets are not displayed).
Observe that in order to fix $A^{(1)}$, we must express it in terms of $F^{(1)}$. To that end, we fix an ansatz: \smash{$A^{(1)}_{\alpha \beta \gamma} = a F^{(1)}_{\alpha \beta \gamma} + b F^{(1)}_{\gamma \alpha \beta} + c F^{(1)}_{\beta \gamma \alpha}$} and then solve explicitly; these are the only three terms~possible~as \smash{$F^{(1)}_{\alpha \beta \gamma} = F^{(1)}_{(\alpha \beta) \gamma}$}. We find that
\[
A^{(1)}_{\alpha \beta \gamma} = -\frac{1}{4}\bigl(F^{(1)}_{\alpha \beta \gamma} + F^{(1)}_{\gamma \alpha \beta} - F^{(1)}_{\beta \gamma \alpha}\bigr)
\]
sets the first order term in equation~\eqref{1-order-correction} to zero identically, and thus we have found $s_{\alpha}$ such that
\[
G_{\alpha \beta} = \delta_{\alpha \beta} + \mathcal{O}\bigl(s^2\bigr) ,
\]
fixing $s_{\alpha}$ to second order in a way dependent on the original choice of extension of $F^{(1)}$.


With this result, we find a useful identity for the normal fundamental form $\beta$ associated with the parallelization of $N \Lambda$:
\begin{align}
n^b_{\alpha} \nabla_a n_{b \beta} \={}& n^b_{[\alpha} \nabla_a n_{b \beta]} + n^b_{(\alpha} \nabla_a n_{b \beta)}
\= n^b_{[\alpha} \bigl(\nabla_a^\top + n_{a \gamma} n^c_{\gamma} \nabla_c \bigr) n_{b \beta]} + \frac{1}{2} \nabla_a n^b_{(\alpha} n_{b \beta)} \nonumber \\
\={}& \beta_{a \alpha \beta} + n_{a \gamma} n^c_{\gamma} n^b_{[\alpha} \nabla_c n_{b \beta]}
\= \beta_{a \alpha \beta} + n_{a \gamma} n^c_{\gamma} n^b_{[\alpha} \nabla_b n_{c \beta]} \nonumber \\
\={}& \beta_{a \alpha \beta} - n_{a \gamma} n^b_{[\alpha} n^c_{\beta]} \nabla_b n_{c \gamma}
\= \beta_{a \alpha \beta} . \label{ndn-Riem}
\end{align}
In the above computation, we used that when $n_{\alpha} = {\rm d} s_{\alpha}$, we evidently have that $\nabla_{[a} n_{b]\beta} = 0$. Observe from this calculation that \smash{$[n_{\alpha},n_{\beta}]^a \= 2 \beta^a_{\alpha \beta}$}. Further, it thus follows that
\begin{align} \label{nabla-n}
\nabla_a n_{b \beta} \= \II_{ab \beta} + n_{b \alpha} \beta_{a \alpha \beta} + n_{a \alpha} \beta_{b \alpha \beta} .
\end{align}


\subsection{Second order} %\label{2o-Riem}
Section~\ref{fo-Riem} established that for a given parallelized Riemannian submanifold embedding, there exists a unique family of defining maps related by $s'_{\alpha} - s_{\alpha} = \mathcal{O}\bigl(s^3\bigr)$ such that
\begin{align} \label{1-order}
G_{\alpha \beta} = \delta_{\alpha \beta} + F^{(2)}_{\alpha \beta \gamma_1 \gamma_2} s_{\gamma_1} s_{\gamma_2} .
\end{align}
As in that section, our goal is to pare down this family by controlling $F^{(2)}$. From equation~(\ref{1-order-correction}), $F^{(2)}$ is uniquely determined in a collar neighborhood of $\Lambda$ given the initial choice of $F^{(1)}$. Our goal will be to adjust $s_{\alpha}$ so that we may remove as much of $F^{(2)}$ as possible.

Using the same reasoning as in Section~\ref{fo-Riem}, we have that $F^{(2)} \in C^\infty (M, \raisebox{3pt}{\scalebox{0.2}{\ydiagram{2,2}}} \; \oplus \raisebox{3pt}{\scalebox{0.2}{\ydiagram{3,1}}} \; \oplus \; \raisebox{1.5pt}{\scalebox{0.2}{\ydiagram{4}}})$. So, we attempt to fix the family of defining functions to cancel $F^{(2)}$ by applying the map
\[
s_{\alpha} \mapsto \tilde{s}_{\alpha} := s_{\alpha} + A^{(2)}_{\alpha \gamma_1 \gamma_2 \gamma_3} s_{\gamma_1} s_{\gamma_2} s_{\gamma_3} ,
\]
which clearly preserves the structure of the Gram matrix condition~(\ref{1-order}) while possibly changing~$F^{(2)}$. However, note that $A^{(2)} \in C^\infty (M, \raisebox{3pt}{\scalebox{0.2}{\ydiagram{3,1}}} \; \oplus \; \raisebox{1.5pt}{\scalebox{0.2}{\ydiagram{4}}})$. Because $A^{(2)}$ and $F^{(2)}$ can live in different representations of $S_k$, it is not always possible to construct $A^{(2)}$ such that $F^{(2)} \mapsto 0$. Nonetheless, one can always find $A^{(2)}$ so that
\[
F^{(2)}_{\alpha \beta \gamma_1 \gamma_2} \mapsto \tilde{F}^{(2)}_{(\alpha \beta)(\gamma_1 \gamma_2)} \in C^\infty (M,\raisebox{3pt}{\scalebox{0.2}{\ydiagram{2,2}}}).
\]



From the above construction, for a given parallelization of $N \Lambda$, there always exists at least one defining map such that
\begin{align*} %\label{2-obs}
G_{\alpha \beta} = \delta_{\alpha \beta} + F^{(2)}_{\alpha \beta \gamma_1 \gamma_2} s_{\gamma_1} s_{\gamma_2} ,
\end{align*}
where $F^{(2)} \in C^\infty (M,\raisebox{3pt}{\scalebox{0.2}{\ydiagram{2,2}}})$ is entirely determined by the parallelization of $N \Lambda$ and the choice of~$F^{(1)}$. We call defining maps that are in this infinite family \textit{associated defining maps}. Along~$\Lambda$, $F^{(2)}|_{\Lambda}$ is a first obstruction to producing a holonomic basis for the extension of the normal bundle. In fact, $F^{(2)}|_{\Lambda}$ can be computed in terms of tensors specified by the parallelization. We now carry out this computation, and going forward we will work with the family of associated defining maps to the parallelization.


Now, observe that
\[
P_{\scalebox{0.2}{\ydiagram{2,2}}}\, n^a_{\gamma_1} n^b_{\gamma_2} \nabla_a \nabla_b G_{\alpha \beta} \= 2 F^{(2)}_{\alpha \beta \gamma_1 \gamma_2} ,
\]
where $P_{\scalebox{0.2}{\ydiagram{2,2}}}$ is the projector onto the $\raisebox{3pt}{{\scalebox{0.2}{\ydiagram{2,2}}}}$ component of a rank-4 tensor in $S^k$. In calculations that will follow, it is useful to have explicit formulas for relevant projection operators. To that end, we find that
\begin{align*}
P_{\scalebox{0.2}{\ydiagram{2,2}}}\, F_{\alpha \beta \gamma \delta} = \frac{1}{12} \bigl(&F_{\alpha \beta \gamma \delta} - F_{\gamma \beta \alpha \delta} - F_{\alpha \delta \gamma \beta} + F_{\gamma \delta \alpha \beta} + F_{\beta \alpha \gamma \delta} - F_{\gamma \alpha \beta \delta} - F_{\beta \delta \gamma \alpha} + F_{\gamma \delta \beta \alpha} \\
&+ F_{\alpha \beta \delta \gamma}- F_{\delta \beta \alpha \gamma} - F_{\alpha \gamma \delta \beta} + F_{\delta \gamma \alpha \beta} + F_{\beta \alpha \delta \gamma} - F_{\delta \alpha \beta \gamma} - F_{\beta \gamma \delta \alpha} + F_{\delta \gamma \beta \alpha}\bigr) ,
\end{align*}
and
\begin{align*}
P_{\scalebox{0.2}{\ydiagram{3,1}}}\, F_{\alpha \beta \gamma \delta} = \frac{1}{8}\bigl(&F_{\alpha \beta \gamma \delta} - F_{\delta \beta \gamma \alpha} + F_{\beta \alpha \gamma \delta} - F_{\delta \alpha \gamma \beta} + F_{\gamma \beta \alpha \delta} - F_{\delta \beta \alpha \gamma} + F_{\alpha \gamma \beta \delta} - F_{\delta \gamma \beta \alpha}
\\&+ F_{\beta \gamma \alpha \delta}- F_{\delta \gamma \alpha \beta} + F_{\gamma \alpha \beta \delta} - F_{\delta \alpha \beta \gamma} \bigr) .
\end{align*}
We can now compute directly
\begin{align*}
P_{\scalebox{0.2}{\ydiagram{2,2}}}\, n^a_{\gamma_1} n^b_{\gamma_2} \nabla_a \nabla_b n^c_{\alpha} n_{c \beta} ={}& P_{\scalebox{0.2}{\ydiagram{2,2}}}\, n^{(a}_{\gamma_1} n^{b)}_{\gamma_2} \bigl[ 2 (\nabla n_{\alpha})^c_{(a} (\nabla n_{\beta})_{b) c} + 2 n^c_{(\alpha|} \nabla_a \nabla_b n_{c| \beta)} \bigr] \\
\={}& -2 \beta_{c \alpha (\gamma_1 } \beta^c_{\gamma_2) \beta} - \frac{2}{3} R_{\gamma_1 (\alpha \beta) \gamma_2} .
\end{align*}
Putting this together, we have that
\begin{align} \label{Riem-2o-obstruction}
F^{(2)}_{\alpha \beta \gamma_1 \gamma_2} \= - \beta_{c \alpha (\gamma_1 } \beta^c_{\gamma_2) \beta} - \frac{1}{3} R_{\gamma_1 (\alpha \beta) \gamma_2} .
\end{align}
It is important to note that $F^{(2)}$ is (as yet) dependent on the choice of extension of $F^{(1)}|_{\Lambda}$. However, what this computation verifies is that this choice of extension yields a unique obstruction $F^{(2)}|_{\Lambda}$ that depends only on the parallelization of $N \Lambda$ and not on some choice of how the parallelization or $F^{(1)}|_{\Lambda}$ extends off $\Lambda$.

Thus, the obstruction in equation~(\ref{Riem-2o-obstruction}) is a true invariant of the parallelized Riemannian submanifold embedding. As such, it may be of independent interest to study those structures for which this obstruction vanishes. While there is indeed an obstruction, we may proceed regardless.




% Original source lines 505--873
\section{Conformal geometry}\label{conf-geom-sec}
We now take a detour to explain the conformal calculus that will be used in the following sections. Much of what is summarized here can be found in more detail in any of~\cite{BEG,CapGover1999,CapGover2003,GOpet,GW}.

Consider a conformal manifold $\bigl(M^d,\cc\bigr)$ with $d \geq 3$. We may view this conformal manifold as a smooth ray subbundle $\mathcal{Q} \subset \odot^2 T^* M$ coordinatized by $\bigl(p,t^2 g\bigr)$ for $p \in M$, $t > 0$, and $g \in \cc$ a~metric on $M$. Indeed, this ray subbundle is a principal $\mathbb{R}_+$ bundle and hence, for each $w \in \mathbb{R}$, induces a natural associated line bundle, called a \textit{density bundle of weight $w$}, associated to $\mathcal{Q}$ via the irreducible representation $\mathbb{R}_+ \ni t \mapsto t^{-w} \in \operatorname{End}(\mathbb{R})$. Such a density bundle is denoted~$\ce M[w]$. Sections of this bundle can then be viewed as functions $\varphi$ on $\mathcal{Q}$ with the property that
\[\varphi\bigl(p, \Omega^2 g\bigr) = \Omega^w \varphi(p,g) .\]
We call such a function a (\textit{conformal}) \textit{density}. For practical purposes, it is often more convenient to view a density as a double equivalence class of functions, $\varphi = [g;f] = \bigl[\Omega^2 g; \Omega^w f\bigr]$, where $f$ is some function in $C^\infty M$. Note that we can tensor a density bundle of weight $w$ with any bundle~$\mathcal{V} M$ over $M$ to obtain a weighted bundle denoted by $\mathcal{V} M[w]$.

Of particular importance is the so-called \textit{conformal metric} $\bg$ which is the tautological section of $\odot^2 T^* M[2]$ given~by $\bg = [g;g]$. In the conformal setting, the conformal metric can be used for raising and lowering indices on tensor-valued densities.

A feature of conformal geometry is that, on the tangent bundle $TM$, there is no connection that respects the conformal structure. Indeed, even on densities, the naive exterior derivative only respects the conformal structure on weight-0 densities. On the other hand, for any $g \in \cc$, there exists a \textit{true scale} $0<\tau \in \Gamma(\ce M[1])$ which is a weight 1 positive density that determines~$g$ via~$g = \bg/\tau^2$. Then, for any true scale $\tau$, we can define a connection on weight-$w$ densities via $\nabla^\tau := \tau^w \circ d \circ \tau^{-w}$. A straightforward calculation shows that for any two true scales $\tau$ and $\tau' = \Omega \tau$ and any $\varphi \in \Gamma(\ce M[w])$, we have that
\[\bigl(\nabla^{\tau'} - \nabla^{\tau}\bigr) \varphi = -w \Upsilon \varphi ,\]
where $\Upsilon := d \log \Omega$. Nonetheless, this family of connections is insufficient for generating invariants of a conformal structure, so instead we turn to tractors. The \textit{standard tractor bundle} over $(M,\cc)$ is a rank $d+2$ vector bundle $\mathcal{T} M$, which for any choice of metric $g \in \cc$ is canonically isomorphic to
\[\mathcal{T}M \stackrel{g}{\cong} \ce M[1] \oplus TM[-1] \oplus \ce M[-1] .\]
We call this isomorphism a \textit{choice of splitting}. In a choice of splitting, we may thus decompose a section of the standard tractor bundle $T \in \Gamma(\ct M)$ according to
\[T^A \stackrel{g}{=} \bigl(\tau^+, \tau^a, \tau^-\bigr) .\]
Any choice of splitting can be related to another according to
\[T^A \stackrel{\Omega^2 g}{=} \biggl(\tau^+, \tau^a + \Upsilon^a \tau^+, \tau^- - \Upsilon_a \tau^a - \frac{1}{2} \Upsilon_a \Upsilon^a \tau^+\biggr) ,\]
where $\Upsilon^a := \bg^{ab} \Upsilon_a$. Observe that $\tau^+$ is independent of choice of splitting and hence is conformally invariant.

Note that for any $T,U \in \Gamma(\ct M)$, there exists an inner product that is independent of the choice of splitting given~by
(in a choice of splitting)
\[T \csdot U \stackrel{g}{=} \bg_{ab} \tau^a \mu^b + \tau^+ \mu^- + \tau^- \mu^+ ,\]
where $U^A \stackrel{g}{=} \bigl(\mu^+, \mu^a, \mu^-\bigr)$. Indeed, this inner product defines the \textit{tractor metric} $h \in \Gamma\bigl( \odot^2 \ct^* M\bigr)$ expressible in a choice of splitting as
\[h_{AB} \stackrel{g}{=} \begin{pmatrix}
0 & 0 & 1 \\
0 & \bg_{ab} & 0 \\
1 & 0 & 0
\end{pmatrix} .\]
Furthermore, this tractor metric provides an isomorphism between the standard tractor bundle and its dual
\[\ct^* M \stackrel{g}{\cong} \ce M[1] \oplus T^* M[1] \oplus \ce M[-1] .\]

The standard tractor bundle also appears in the construction of other useful bundles. By~tensoring the standard tractor bundle with the weight-$1$ density bundle, we obtain $\ct M[1]$, which has a canonically defined section called the \textit{canonical tractor}, given in any choice of splitting by $X^A \stackrel{g}{=} (0,0,1) \in \Gamma(\ct M[1])$. More generally, we shall use the terminology \textit{tractor bundle} to refer to any tensor product of copies of the standard tractor bundle and density bundles.

Notably, there exists a well-defined and canonical \textit{standard tractor connection} on $\mathcal{T}M$,
\begin{align*}
\nabla^\ct \colon \ \Gamma(\ct M) &\rightarrow \Gamma(T^* M \otimes \ct M) \\
T^B &\mapsto \nabla^\ct_a T^B \stackrel{g}{=} \begin{pmatrix} \nabla_a \tau^+ - \tau_a \\ \nabla_a \tau^b + \bg_a^b \tau^- + (P^g)_a^b \tau^+ \\ \nabla_a \tau^- - P^g_{ab} \tau^b \end{pmatrix} .
\end{align*}
It is this connection that ensures that the 1-form-valued tractor above is independent of the choice of splitting (and hence is a section of a tractor bundle). Note that this connection can be extended to tractor bundle. However, this connection does \textit{not} map to $\mathcal{T}^* M \times \mathcal{T}M$. Instead, we can use the tractor connection to build a tractor-valued operator known as the \textit{Thomas-$D$ operator}, given in a choice of splitting by
\begin{align*}
D_A \colon \ \Gamma\bigl(\ct^\Phi M[w]\bigr) &\rightarrow \Gamma\bigl(\ct^* M \otimes \ct^\Phi M[w-1]\bigr), \\
T^\Phi &\mapsto D_A T^\Phi \stackrel{g}{=} \begin{pmatrix}
(d+2w-2)w T^\Phi \\
(d+2w-2) \nabla^\ct_a T^\Phi \\
-(\Delta^\ct + w J^g) T^\Phi
\end{pmatrix} ,
\end{align*}
where $\Delta^\ct = \bg^{ab} \nabla^\ct_a \nabla^\ct_b$ is the tractor Laplacian. From the Thomas-$D$ operator, we can define another useful operator on $\ct^\Phi M[w]$ -- the \textit{hatted} Thomas-$D$ operator
\[\hd_A = \frac{1}{d+2w-2} D_A ,\]
which is well defined so long as $w \neq 1 - \frac{d}{2}$. While the Thomas-$D$ operator is not a derivation, it has the useful feature that $D_A \circ D^A = 0$. Furthermore, its failure to be a derivation is well controlled and given~by the following formula~\cite{Will1}:
\begin{gather*}
\hd^A (T_1 \cdots T_\ell) - \sum_{i=1}^\ell T_1 \cdots \bigl(\hd^A T_i\bigr) \cdots T_\ell \\
\qquad{}= - \frac{2 X^A}{d+2 \sum_{i=1}^\ell w_i -2} \sum_{1 \leq i < j \leq \ell} T_1 \cdots \bigl(\hd^B T_i\bigr) \cdots \bigl(\hd_B T_j\bigr) \cdots T_\ell ,
\end{gather*}
where for each $1 \leq i \leq \ell$, \smash{$T_i \in \Gamma\bigl(\ct^\Phi M[w_i]\bigr)$}, and the weight constraints $d+2\sum_{i=1}^{\ell} w_i - 2 \neq 0$ and \smash{$\big\{w_i \neq 1 - \frac{d}{2}\big\}_{i=1}^{\ell}$} are satisfied.
While it is clear that this operator is not a covariant derivative, it often plays the role of one. Indeed, the commutator of two hatted Thomas-$D$ operators produces a curvature-like tractor called the $W$-tractor, according to
\[\bigl[\hd_A, \hd_B\bigr] T^D = W_{AB}{}^D{}_E T^E + \frac{4}{d+2w-4} X_{[A} W_{B] C}{}^D{}_E \hd^C T^E ,\]
where such a commutator is defined. The $W$-tractor has Riemann-like symmetries, is trace-free, and satisfies
\[X^A W_{ABCD} = 0 = \hd^A W_{ABCD} .\]
This tractor is composed of the Weyl tensor, the Cotton tensor, and a generalized notion of the Bach tensor. See~\cite{BGW1, GOmin} for more details. Note that we will use the same letter for both the $W$-tractor and the Weyl tensor, but it will be clear from context (i.e., which indices are being used) which object is appearing in a given formula.

Because all of the aforementioned tractor objects respect the conformal structure of $(M,\cc)$, any combination of these tractors and tractor-valued operators can be used to produce a conformally invariant quantity. It is this feature that makes the tractor calculus analogous to the Ricci calculus of Riemannian geometry, and we will utilize this advantage heavily in the sections that follow.

\section[Orthonormal frames for normal bundles of conformal submanifold embeddings]{Orthonormal frames for normal bundles\\ of conformal submanifold embeddings} \label{conf-frames}
We now turn to the case of conformal submanifold embedings $\Lambda^{d-k} \hookrightarrow \bigl(M^d,\cc\bigr)$. As the tangent bundle decomposition for a Riemannian submanifold embedding in equation~(\ref{TM-decomp}) only depends on the orthogonality of vectors rather than their lengths, a conformal structure admits an identical decomposition. Furthermore, given any representative $g \in \cc$ and an orthonormal frame $\{(n^g)^a_{\alpha}\}$ with respect to $g$ of $N \Lambda$, we can promote such an orthonormal frame to an orthonormal frame-valued density (OFVD) $\{\bn^a_{\alpha}\} := \{[g ; (n^g)^a_{\alpha}]\}$ with weight $-1$. Then it follows that $\bg(\bn_{\alpha}, \bn_{\beta}) = \delta_{\alpha \beta}$.

Interestingly, normal fundamental forms are conformally-invariant. Indeed, given an OFVD and a metric representative $g \in \cc$, by direct computation of the conformal transformation of the Levi-Civita connection and the orthonormal frame-valued density, it follows that $\beta_{a \alpha \beta} \in \Gamma(T^* \Sigma[0] \times \raisebox{3pt}{\scalebox{0.2}{\ydiagram{1,1}}})$. So, the normal curvature is also invariant. Thus all of the results of Section~\ref{riem-frames} directly generalize to the equivalent conformal structures without additional consideration. Unfortunately, the results of Section~\ref{unit-def-funcs} do not generalize directly in the conformal setting. Instead, additional obstructions and other troublesome features arise.

\begin{Remark}
The conformal invariance of the normal fundamental form may seem to contradict the observation that, for a curve $\Lambda^1 \hookrightarrow \mathbb{R}^3$, the normal fundamental form of the Frenet frame picks out the torsion. However, this observation is naive. The normal fundamental form is conformally invariant when the orthonormal frame vectors are density valued. However, rescaling of the metric generically changes the plane of osculation for a given curve and hence rotates the Frenet frame. Thus the failure of a spacecurve's torsion to be conformally-invariant is a~consequence of the changing orientation of the Frenet frame.
\end{Remark}

\section[Defining densities for parallelized conformal submanifold embeddings]{Defining densities for parallelized conformal\\ submanifold embeddings} \label{conf-densities-sec}

As in the Riemannian setting, we would like to better study the conformal submanifold embedding $\Lambda^{d-k} \hookrightarrow \bigl(M^d,\cc\bigr)$ by studying the jets of geometrically-determined \textit{defining densities} -- densities of weight 1 that, in any metric representative, are represented by a defining map for $\Lambda$. To do so, we will attempt to construct the conformal analog of associated defining maps. We begin by picking any $g \in \cc$ and constructing the canonical defining map $s_{\alpha}^g$ for the provided OFVD in that metric representative. We can then promote this associated defining map to a defining density $\sigma_{\alpha}$ of weight 1. (Note that we are abusing language by calling an element of $\Gamma(\ce M[1] \times \raisebox{2pt}{\scalebox{0.2}{\ydiagram{1}}})$ a defining density.) Had we picked another metric conformally related to the first, the difference between the two densities would be at least second order in a defining density. For that reason, the initial choice of arbitrary metric to begin the construction is irrelevant for what follows (as we will change higher order terms anyway). That is, for a given OFVD, modulo terms of second order in a defining density, there exists a unique defining density $\sigma_{\alpha} = \bigl[g; s^g_{\alpha}\bigr] \in \Gamma(\ce M[1])$ such that \smash{$\mathrm{d} \sigma_{\alpha} \= \bg(\bn_{\alpha},\cdot)$}. (Note that $\mathrm{d} \sigma_{\alpha}$ can be defined given a choice of scale and depends on this choice. However, it is easy to see that its restriction to $\Lambda$ is independent of this choice, and hence can be regarded as a section of $T^*M[1]|_{\Lambda}$.) Our goal will be to maximally improve this density so that we can extract invariants of the parallelized conformal submanifold embedding.

In general, the Gram matrix of a defining map is not conformally invariant away from $\Lambda$. Even in the codimension 1 case, for two metric representatives $g, \Omega^2 g \in \cc$, we find that
\[\bigl|{\rm d} s^{\Omega^2 g}\bigr|^2_{\Omega^2 g} = |{\rm d} (\Omega s^g)|^2_{\Omega^2 g} = |\Omega (n^g_a + s \Upsilon_a)|^2_{\Omega^2 g} = |n^g_a + s^g \Upsilon_a|^2_g \neq |{\rm d}s^g|^2_g .\]
From this computation, we have shown explicitly that $G^{\Omega^2 g} \neq G^g$ except on the submanifold $\Lambda$. Going forward, we will be dropping the superscripts denoting metric representatives when the context is clear.

The natural analog of the (Riemannian) Gram matrix for a defining map to the conformal case is the \textit{conformal Gram matrix} of a defining density. Given a defining density for a~submanifold~$\sigma_\alpha$, we can build a family of \textit{scale tractors},
\[\big\{N^A_{\alpha} := \hd^A \sigma_{\alpha} \big\}_{\alpha = 1}^k .\]
The conformal Gram matrix is defined according to
\[\bG_{\alpha \beta} := h_{AB} N^A_{\alpha} N^B_{\beta} .\]
The conformal Gram matrix is conformally invariant by construction. Furthermore, upon restriction to the submanifold $\Lambda$, in any representative $g \in \cc$ the conformal Gram matrix agrees with the Riemannian Gram matrix of the representative of $\sigma_{\alpha}$ given~by $[g; s_{\alpha}]$. Indeed, for $g \in \cc$ any metric representative and $\sigma_{\alpha} = [g;s_{\alpha}]$ a defining density for $\Lambda^{d-k} \hookrightarrow (M,\cc)$, the components of $\bG_{\alpha \beta}$ are given~by
\[\bG_{\alpha \beta} = [g; G_{\alpha \beta}] := \biggl[g; (\nabla^a s_{\alpha})(\nabla_a s_{\beta}) - \frac{2}{d} s_{(\alpha} (\Delta^g + J^g) s_{\beta)}\biggr] .\]
Note that we are overloading notation: going forward we use $G_{\alpha \beta}$ to refer to a representative of the conformal Gram matrix, rather than the Riemannian Gram matrix. However note that every representative is the same, as the conformal Gram matrix has weight 0 -- we merely mean that we refer not to the double equivalence class $[g;G_{\alpha \beta}]$ but just $G_{\alpha \beta}$.



In the Riemannian setting, we found that, for a given parallelization of $N \Lambda$, a unique family of canonical defining maps could be obtained to arbitrary order. Further, when the background was flat and the normal curvature vanished, we found that this family consists entirely of unit defining maps. We proceed now with a similar program in the conformal setting.

\subsection{First order} \label{conf-first-order}
Given an initial choice of defining density built from an associated defining map for $\Lambda \hookrightarrow (M,\cc)$, it is easy to see that any representative of the conformal Gram matrix agrees with the Riemannian Gram matrix for the same associated defining map along $\Lambda$. It follows then that
\[\bG_{\alpha \beta} = \delta_{\alpha \beta} + F^{(1)}_{\alpha \beta \gamma} \sigma_{\gamma} \]
for some $F^{(1)} \in \Gamma(\ce M[-1] \times (\raisebox{3pt}{\scalebox{0.2}{\ydiagram{2,1}}} \; \oplus \; \raisebox{2pt}{\scalebox{0.2}{\ydiagram{3}}}))$, similar to the Riemannian setting. As in that setting, we will pick \textit{some} $F^{(1)}$ that extends $F^{(1)}|_{\Lambda}$, with the goal of showing that this choice is immaterial. Note, however, both the trace and trace-free parts of these representations will play a role in what follows.

We will now improve $\sigma$ at second order to cancel $F^{(1)}$ along $\Lambda$. Let
\[\tilde{\sigma}_{\alpha} = \sigma_{\alpha} + A^{(1)}_{\alpha \gamma_1 \gamma_2} \sigma_{\gamma_1} \sigma_{\gamma_2} ,\]
where $A^{(1)} \in \Gamma(\ce M[-1] \times (\raisebox{3pt}{\scalebox{0.2}{\ydiagram{2,1}}} \; \oplus \; \raisebox{2pt}{\scalebox{0.2}{\ydiagram{3}}}))$. Going forward, we will implicitly work in a fixed scale $g \in \cc$. To begin, we compute
\begin{align*}
&\tilde{n}_{\alpha} = n_{\alpha} + 2 A^{(1)}_{\alpha(\gamma_1 \gamma_2)} n_{\gamma_1} s_{\gamma_2} + \mathcal{O}\bigl(s^2\bigr), \\
& \Delta \tilde{s}_{\alpha} = \Delta s_{\alpha} + 2 A^{(1)}_{\alpha \gamma \gamma} + \mathcal{O}(s) .
\end{align*}
Thus, we have that
\begin{align*}
\tilde{G}_{\alpha \beta} &{}= \tilde{n}^a_{\alpha} \tilde{n}_{a \beta} - \frac{2}{d} \tilde{s}_{(\alpha} (\Delta^g + J^g) \tilde{s}_{\beta)} \\
&{}= n^a_{\alpha} n_{a \beta} + 4 n^a_{\gamma_1} n_{a (\alpha} A^{(1)}_{\beta)(\gamma_1 \gamma_2)} s_{\gamma_2} - \frac{2}{d} s_{(\alpha}(\Delta^g + J^g)s_{\beta)} - \frac{4}{d} s_{(\alpha} A^{(1)}_{\beta) \gamma \gamma} + \mathcal{O}\bigl(s^2\bigr) \\
&{}= G_{\alpha \beta} + 4 A^{(1)}_{(\alpha \beta) \gamma} s_{\gamma} - \frac{4}{d} s_{(\alpha} A^{(1)}_{\beta) \gamma \gamma} + \mathcal{O}\bigl(s^2\bigr)\\
&{}= \delta_{\alpha \beta} + \biggl(F^{(1)}_{\alpha \beta \omega} + 4 A^{(1)}_{(\alpha \beta) \omega} - \frac{4}{d} \delta_{\omega (\alpha} A^{(1)}_{\beta) \gamma \gamma}\biggr) s_{\omega} + \mathcal{O}\bigl(s^2\bigr) .
\end{align*}
Indeed, we see that a trace of $A^{(1)}$ appears here. Thus we must be careful to check that all irreducible components of $F^{(1)}$ can be cancelled. In this case, we have that
\[F^{(1)} = \bigl(\raisebox{3pt}{\scalebox{0.4}{\ydiagram{2,1}}}_{\,\circ} \oplus \raisebox{0pt}{\scalebox{0.4}{\ydiagram{1}}}\bigr) \; \oplus \; (\scalebox{0.4}{\ydiagram{3}}_{\,\circ} \oplus \raisebox{0pt}{\scalebox{0.4}{\ydiagram{1}}}) ,\]
and similarly for $A^{(1)}$. Note that the trace-free parts of $F^{(1)}$ can be cancelled for the same reason they can be cancelled in the Riemannian case. To force the traceful parts to vanish, we demand that
\begin{align*}
&0= F^{(1)}_{\alpha \alpha \omega} + 4 A^{(1)}_{\alpha \alpha \omega} - \frac{4}{d} A^{(1)}_{\omega \alpha \alpha} \;\;\text{ and } \\
&0= F^{(1)}_{\omega \alpha \alpha} + 2 A^{(1)}_{\omega \alpha \alpha} + 2 A^{(1)}_{\alpha \alpha \omega} - \frac{2}{d} A^{(1)}_{\omega \alpha \alpha} - \frac{2k}{d} A^{(1)}_{\omega \alpha \alpha} \\
& \hphantom{0}{}= F^{(1)}_{\omega \alpha \alpha} + 2 A^{(1)}_{\alpha \alpha \omega} + \frac{2(d-k-1)}{d} A^{(1)}_{\omega \alpha \alpha} .
\end{align*}
The determinant of this linear system \smash{\big(with $A^{(1)}_{\alpha \alpha \omega}$ and $A^{(1)}_{\omega \alpha \alpha}$ as variables\big)} is $8(d-k)/d$, and thus for all $k \neq d$, we can solve for the traces of $A^{(1)}$ such that they cancel the traces of $F^{(1)}$. We have thus shown the intermediate proposition:
\begin{Proposition}
Let $\Lambda^{d-k} \hookrightarrow \bigl(M^d,\cc\bigr)$ be a conformal submanifold embedding with a parallelization given in a choice of metric representative by $\{n_{\alpha}\}$. Then there exists a defining map $\sigma_{\alpha} \in \Gamma(\ce M[1] \times \raisebox{2pt}{\scalebox{0.2}{\ydiagram{1}}})$ unique up to order $\sigma^3$ such that $\mathrm{d} \sigma_{\alpha}|_{\Lambda} = [g; n_{a \alpha}]$ and its conformal Gram matrix satisfies
\[\bG_{\alpha \beta} = \delta_{\alpha \beta} + \mathcal{O}\bigl(\sigma^2\bigr) .\]
\end{Proposition}

Given such a defining map that satisfies $\bG_{\alpha \beta} = \delta_{\alpha \beta} + \mathcal{O}\bigl(\sigma^2\bigr)$, it will be useful to evaluate~$N^A_{\alpha}|_{\Lambda}$. Doing so requires the evaluation of $\nabla \csdot n_{\alpha}|_{\Lambda}$ in a choice of metric representative. First, note that, unlike in the Riemannian case, we have that
\[n^a_{\alpha} n_{a \beta} = \delta_{\alpha \beta} + \frac{2}{d} s_{(\alpha} \Delta s_{\beta)} + \mathcal{O}\bigl(s^2\bigr) .\]
Hence, following (with a slight modification) the computation leading to equation~(\ref{ndn-Riem}), we find that
\[n^b_{\alpha} \nabla_a n_{b \beta} \= \beta_{a \alpha \beta} + \frac{1}{d} n_{a \alpha} \nabla \csdot n_{\beta} .\]
Consequently, we have that
\begin{align*}% \label{dn-conf}
\nabla_a n_{b \beta} \= \II_{ab \beta} + n_{b \alpha} \beta_{a \alpha \beta} + n_{a \alpha} \beta_{b \alpha \beta} + \frac{1}{d}n_{a \alpha} n_{b \alpha} \nabla \csdot n_{\beta} .
\end{align*}
It thus follows that
\[\nabla \csdot n_{\beta} \=  d  H_{\beta} ,\]
where $H_{\beta} = \frac{1}{d-k} \II^a_{a\beta}$ is the mean curvature, and so
\[N^A_{\alpha} \stackrel{\Lambda,g}{=} \begin{pmatrix}
0\\ \, n^a_{\alpha}
\\
-H_{\alpha}
\end{pmatrix} .\]
We may also compute the tractor equivalent of the bulk formula for the normal fundamental form, given~by
\[B^A_{\alpha \beta} := N_{[\alpha} \csdot \hd N^A_{\beta]} \in \Gamma(\ct M[-1] \times \raisebox{2.5pt}{\scalebox{0.2}{\ydiagram{1,1}}} ) .\]
By direct computation, this tractor has the property that
\[X_A B^A_{\alpha \beta} = \hd_A B^A_{\alpha \beta} = 0 .\]
This tractor necessarily takes the form
\[B^A_{\alpha \beta} \stackrel{g}{=} \begin{pmatrix}
0\\
\, n^b_{[\alpha} \nabla_b n^a_{\beta]} - \frac{1}{d} n^a_{[\alpha} \nabla \csdot n_{\beta]} + \mathcal{O}(s)\\
\ast
\end{pmatrix}
\stackrel{\Lambda,g}{=}
\begin{pmatrix}
0\\
\, \beta^a_{\alpha \beta} \\
\ast
\end{pmatrix}
 .\]
Finally, it is useful to express the projecting part of the tractor equivalent of $\nabla_a n_b$,
\begin{align*}
P_{AB \alpha} :={}& \hd_A N_{B \alpha}\\
 ={}& \begin{pmatrix}
0 & 0 & 0 \\
0 & \nabla_a n_{b \alpha} + s_{\alpha}P_{ab} + g_{ab} \rho_{\alpha} & \ast \\
0 & \ast & \ast
\end{pmatrix}
 \= \begin{pmatrix}
0 & 0 & 0 \\
0 & \IIo_{ab \alpha} + 2 n_{(a \beta} \beta_{b) \beta \alpha} & \ast \\
0 & \ast & \ast
\end{pmatrix} ,
\end{align*}
where $\rho_{\alpha} := -\frac{1}{d}(\Delta s_{\alpha} + J s_{\alpha})$ and \smash{$\IIo_{\alpha} := \II_{\alpha} - \bar{g} H_{\alpha} \in \Gamma\bigl(\odot^2_{\circ} T^* \Lambda[1]\bigr)$} is the trace-free component of the second fundamental form.
With this tractor in hand, we define a scalar-valued density called the \textit{submanifold rigidity density} $K_{\alpha \beta} := P_{AB \alpha} P^{AB}_{\beta}$ which is ubiquitous. A short calculation shows that
\[K_{\alpha \beta} \= \IIo^2_{\alpha \beta} - 2 \beta_{b \gamma (\alpha} \beta^b_{\beta) \gamma} .\]
We now proceed to higher orders.

\subsection{Second order}

As in the Riemannian case, we now proceed to adjust $\sigma_{\alpha}$ at third order to eliminate as much from the second order term $\bG_{\alpha \beta}$ as possible. Pulling from the Riemannian case, we expect to fail here as well. Suppose that $\sigma_{\alpha}$ has
\[\bG_{\alpha \beta} = \delta_{\alpha \beta} + F^{(2)}_{\alpha \beta \gamma_1 \gamma_2} \sigma_{\gamma_1} \sigma_{\gamma_2} .\]
As a reminder, $F^{(2)}$ may depend on the initial choice made of $F^{(1)}$. Now, we can decompose~$F^{(2)}$ into irreducible representations so that
\[
F^{(2)} \in
\raisebox{2.5pt}{\scalebox{0.4}{\ydiagram{2,2}}}_{\,\circ}
 \; \oplus \;
 \raisebox{2.5pt}{\scalebox{0.4}{\ydiagram{3,1}}}_{\,\circ}
 \; \oplus \;
 \raisebox{0pt}{\scalebox{0.4}{\ydiagram{4}}}_{\,\circ}
\; \oplus \;
 3\, \raisebox{1pt}{\scalebox{0.4}{\ydiagram{2}}}_{\,\circ}
\; \oplus \;
 \raisebox{2.5pt}{\scalebox{0.4}{\ydiagram{1,1}}}
 \; \oplus \;
2.
\]
On the other hand, we can define $\tilde{\sigma}_{\alpha} = \sigma_{\alpha} + A^{(2)}_{\alpha \gamma_1 \gamma_2 \gamma_3 } \sigma_{\gamma_1} \sigma_{\gamma_2} \sigma_{\gamma_3}$. Then we can decompose $A^{(2)}$ in a similar fashion, yielding
\[A^{(2)} \in
 \raisebox{2.5pt}{\scalebox{0.4}{\ydiagram{3,1}}}_{\,\circ}
 \; \oplus \;
 \raisebox{0pt}{\scalebox{0.4}{\ydiagram{4}}}_{\,\circ}
\; \oplus \;
 2\, \raisebox{1pt}{\scalebox{0.4}{\ydiagram{2}}}_{\,\circ}
\; \oplus \;
 \raisebox{2.5pt}{\scalebox{0.4}{\ydiagram{1,1}}}
 \; \oplus \;
1.
\]
 So naively one might predict there will be obstructions of the form
\[
F^{(2)} \in
\raisebox{2.5pt}{\scalebox{0.4}{\ydiagram{2,2}}}_{\,\circ}
 \; \oplus \;
 \, \raisebox{1pt}{\scalebox{0.4}{\ydiagram{2}}}_{\,\circ}
 \; \oplus \;
1.
\]
To check this naive understanding, we will explicitly compute for the last time in this case. As before, we have that
\begin{align*}
&\tilde{n}_{\alpha}= n_{\alpha} + 3 A^{(2)}_{\alpha (\gamma_1 \gamma_2 \gamma_3)} s_{\gamma_1} s_{\gamma_2} n_{\gamma_3} + \mathcal{O}\bigl(s^3\bigr), \\
&\Delta \tilde{s}_{\alpha}= \Delta s_{\alpha} + 6 A^{(2)}_{\alpha \gamma_1 \gamma \gamma} s_{\gamma_1} + \mathcal{O}\bigl(s^2\bigr) .
\end{align*}
Plugging these into the formula for $G_{\alpha \beta}$, we have
\begin{align*}
\tilde{G}_{\alpha \beta} &{}= \tilde{n}^a_{\alpha} \tilde{n}_{a \beta} - \frac{2}{d} \tilde{s}_{(\alpha} (\Delta^g + J^g) \tilde{s}_{\beta)} \\
&{}= n^a_{\alpha} n_{a \beta} - \frac{2}{d} s_{(\alpha} (\Delta^g + J^g) s_{\beta)} + \biggl(6 A^{(2)}_{(\alpha \beta) \omega_1 \omega_2} - \frac{12}{d} \delta_{\omega_1 (\alpha} A^{(2)}_{\beta) \gamma \gamma \omega_2}\biggr) s_{\omega_1} s_{\omega_2} \\
&{}= \delta_{\alpha \beta} + \biggl(F^{(2)}_{\alpha \beta \omega_1 \omega_2} + 6 A^{(2)}_{(\alpha \beta) \omega_1 \omega_2} - \frac{12}{d} \delta_{\omega_1 (\alpha} A^{(2)}_{\beta) \gamma \gamma \omega_2}\biggr) s_{\omega_1} s_{\omega_2} .
\end{align*}
Note that there is implicit symmetrization over $\omega_1$ and $\omega_2$ in these displays. Just as in the Riemannian case, it is straightforward to see that the $ \raisebox{2.5pt}{\scalebox{0.2}{\ydiagram{3,1}}}_{\,\circ}$ and $\raisebox{1.5pt}{\scalebox{0.2}{\ydiagram{4}}}_{\,\circ} $ components of $F^{(2)}$ can be cancelled directly, and it is impossible to cancel the $\raisebox{2.5pt}{\scalebox{0.2}{\ydiagram{2,2}}}_{\,\circ}$ component of $F^{(2)}$.
We now consider the remaining components of $F^{(2)}$. If the trace-free parts of this coefficient are to vanish, we must demand that
\begin{align*}
&0= F^{(2)}_{\rho \rho (\omega_1 \omega_2)_\circ} + 6 A^{(2)}_{\rho \rho (\omega_1 \omega_2)_\circ} - \frac{12}{d} A^{(2)}_{(\omega_1 \omega_2)_\circ \rho \rho}, \\
&0
= F^{(2)}_{(\omega_1 \omega_2)_\circ \rho \rho} + \frac{6(d-2)}{d} A^{(2)}_{(\omega_1 \omega_2)_\circ \rho \rho}, \\
&0
= F^{(2)}_{\rho (\omega_1 \omega_2)_\circ \rho} + 3 A^{(2)}_{\rho \rho (\omega_1 \omega_2)_\circ} + \frac{3(d-k-2)}{d} A^{(2)}_{(\omega_1 \omega_2)_\circ \rho \rho}, \\
&0
= F^{(2)}_{\rho [\omega_1 \omega_2] \rho} + \frac{3(d-k-2)}{d} A^{(2)}_{[\omega_1 \omega_2] \rho \rho}, \\
&0
= F^{(2)}_{\rho \rho \gamma \gamma} + \frac{6(d-2)}{d} A^{(2)}_{\rho \rho \gamma \gamma}, \\
&0
= F^{(2)}_{\rho \gamma \gamma \rho} + \frac{6(d-k-1)}{d} A^{(2)}_{\rho \rho \gamma \gamma} .
\end{align*}
To more easily visualize this system of equations, we denote \smash{$x_1 = A^{(2)}_{\rho \rho (\omega_1 \omega_2)_\circ}$, $x_2 = A^{(2)}_{(\omega_1 \omega_2)_\circ \rho \rho}$}, \smash{$x_3 = A^{(2)}_{[\omega_1 \omega_2] \rho \rho}$}, and \smash{$x_4 = A^{(2)}_{\gamma \gamma \rho \rho}$}. Then the matrix describing the linear system of equations is given~by
\[
\begin{pmatrix}
6 & -\frac{12}{d} & 0 & 0 \\[1mm]
0 & \frac{6(d-2)}{d} & 0 & 0 \\[1mm]
3 & \frac{3(d-k-2)}{d} & 0 & 0 \\[1mm]
0 & 0 & \frac{3(d-k-2)}{d} & 0 \\[1mm]
0 & 0 & 0 & \frac{6(d-2)}{d} \\[1mm]
0 & 0 & 0 & \frac{6(d-k-1)}{d}
\end{pmatrix} .
\]
This matrix has rank at most 4. Furthermore, when $k = d-2$, this matrix has rank 3. \big(Note that the rank also drops when $k=1$, but all components of $F^{(2)}$ collapse to a single function, which is uniquely determined by fixing the single function $A^{(2)}$.\big) So as expected, we can remove at most 4 traced degrees of freedom from \smash{$F^{(2)}$} using \smash{$A^{(2)}$}, except when $k = d-2$, in which case we can only remove 3 degrees of freedom. Indeed, for $k \neq d-2$, we have an obstruction given~by
\[
F^{(2)} =
\raisebox{2.5pt}{\scalebox{0.4}{\ydiagram{2,2}}}_{\,\circ}
 \; \oplus \;
 \, \raisebox{1pt}{\scalebox{0.4}{\ydiagram{2}}}_{\,\circ}
 \; \oplus \;
1.
\]
When $k = d-2$, we instead have an obstruction of the form
\[
F^{(2)} =
\raisebox{2.5pt}{\scalebox{0.4}{\ydiagram{2,2}}}_{\,\circ}
 \; \oplus \;
 \, \raisebox{1pt}{\scalebox{0.4}{\ydiagram{2}}}_{\,\circ}
 \; \oplus \;
 \, \raisebox{1pt}{\scalebox{0.4}{\ydiagram{1,1}}}_{\,\circ}
 \; \oplus \;
1.
\]
For the sake of uniformity, when $k \neq d-2$, we will choose to eliminate as much of the components of $F^{(2)}$ in $\raisebox{2.5pt}{\scalebox{0.2}{\ydiagram{3,1}}}$ and $\raisebox{1.5pt}{\scalebox{0.2}{\ydiagram{4}}}$ as possible. Certainly, as in the Riemannian case, we can easily remove the trace-free parts of these components, so we examine the traces.

First, observe that
\[\delta^{\alpha \beta} \delta^{\gamma \delta} F^{(2)}_{(\alpha \beta \gamma \delta)} = \frac{1}{3} \bigl(F^{(2)}_{\alpha \alpha \beta \beta} + 2 F^{(2)}_{\alpha \beta \beta \alpha}\bigr) .\]
So we demand that
\[F^{(2)}_{\alpha \alpha \beta \beta} + 2 F^{(2)}_{\alpha \beta \beta \alpha} + \frac{6(3d-4-2k)}{d} A_{\alpha \alpha \beta \beta}^{(2)} = 0 .\]
This has solutions so long as $(d,k) \neq (2,1)$, and thus we can always remove the full trace of \smash{$F^{(2)}_{(\alpha \beta \gamma \delta)}$} in that setting. In fact, when $(d,k) = (2,1)$, we may directly compute this full trace term (which is in fact the only possibly obstruction, as $k=1$) and see that it vanishes. However, as $A^{(2)}$ is not determined by the total trace of $F^{(2)}$, we may not proceed there.

Next we consider the trace-free symmetric part of the trace of the totally-symmetric component. By the same considerations as before, this component is given~by
\[\frac{1}{6}\bigl(F_{\gamma \gamma (\alpha \beta)_\circ} + 4 F_{\gamma (\alpha \beta)_\circ \gamma} + F_{(\alpha \beta)_\circ \gamma \gamma}\bigr) .\]
Thus, we demand that
\begin{align} \label{tfs-fully-symmetric}
F_{\gamma \gamma (\alpha \beta)_\circ} + 4 F_{\gamma (\alpha \beta)_\circ \gamma} + F_{(\alpha \beta)_\circ \gamma \gamma} + 18 A^{(2)}_{\gamma \gamma (\alpha \beta)_\circ} + \frac{6(3d-2k-8)}{d} A^{(2)}_{(\alpha \beta)_\circ \gamma \gamma} = 0 .
\end{align}
To fully determine both of the components of $A^{(2)}$ above, we can also demand that the trace-free symmetric part of the trace of the $\raisebox{2.5pt}{\scalebox{0.2}{\ydiagram{3,1}}}$ component of $F^{(2)}$ vanishes. This component of $F^{(2)}$ can be expressed as
\[\frac{1}{6}\bigr(F^{(2)}_{(\alpha \beta)_\circ \gamma \gamma} - F^{(2)}_{\gamma \gamma (\alpha \beta)_\circ}\bigl) .\]
So, we must have that
\begin{align} \label{tfs-pistol}
F^{(2)}_{(\alpha \beta)_\circ \gamma \gamma} - F^{(2)}_{\gamma \gamma (\alpha \beta)_\circ} - 6 A^{(2)}_{\gamma \gamma (\alpha \beta)_\circ} + 6 A^{(2)}_{(\alpha \beta)_\circ \gamma \gamma} =0 .
\end{align}
Combining equations~(\ref{tfs-fully-symmetric}) and~(\ref{tfs-pistol}), we find that the determinant of the system is $\frac{72}{d} (3d-k-4)$, which never vanishes for $1 \leq k \leq d-1$. Thus, we can completely remove the trace-free symmetric pieces of both $\raisebox{2.5pt}{\scalebox{0.2}{\ydiagram{3,1}}}$ and $\raisebox{1.5pt}{\scalebox{0.2}{\ydiagram{4}}}$.

Finally, we wish to remove the antisymmetric part of the trace of the $\raisebox{2.5pt}{\scalebox{0.2}{\ydiagram{3,1}}}$ component of $F^{(2)}$, which is given~by \smash{$F^{(2)}_{\gamma [\alpha \beta] \gamma}$}. So long as $k \neq d-2$, this piece can be removed as well. This accounts for all of the degrees of freedom in $A^{(2)}$, and thus we are left with a completely undetermined $\raisebox{2.5pt}{\scalebox{0.2}{\ydiagram{2,2}}}$ component of $F^{(2)}$, precisely as in the Riemannian setting.



Having fixed $A^{(2)}$ and in so doing, forced $F^{(2)} \in \Gamma(\ce M[-2] \times \raisebox{2.5pt}{\scalebox{0.2}{\ydiagram{2,2}}})$, it is now instructive to compute this obstruction.
We begin by noting that
\begin{align} \label{window-projection-bG}
\bG_{\alpha \beta} = \delta_{\alpha \beta} + \bigl(P_{\scalebox{0.2}{\ydiagram{2,2}}}\, F^{(2)}_{\alpha \beta \gamma_1 \gamma_2} \bigr) \sigma_{\gamma_1} \sigma_{\gamma_2} .
\end{align}
Furthermore, we have that
\[ P_{\scalebox{0.2}{\ydiagram{2,2}}}\, F^{(2)}_{\alpha \beta \gamma_1 \gamma_2} = \frac{1}{3} \bigl(2 F^{(2)}_{\alpha \beta \gamma_1 \gamma_2} + 2 F^{(2)}_{\gamma_1 \gamma_2 \alpha \beta} - 2 F^{(2)}_{\gamma_1 (\alpha \beta) \gamma_2} - 2F^{(2)}_{\gamma_2 (\alpha \beta) \gamma_1} \bigr) .\]
Now consider the following projection:
\[P_{\scalebox{0.2}{\ydiagram{2,2}}}\, N^A_{\gamma_1} N^B_{\gamma_2} \hd_A \hd_B \bG_{\alpha \beta} .\]
Using equation~(\ref{window-projection-bG}) and simplifying, we find that
\begin{align}
P_{\scalebox{0.2}{\ydiagram{2,2}}}\, N^A_{\gamma_1} N^B_{\gamma_2} \hd_A \hd_B \bG_{\alpha \beta} &{}\=
2F^{(2)}_{\alpha \beta \gamma_1 \gamma_2}
+ \frac{2}{9(d-2)} \delta_{\alpha \beta} \bigl(F^{(2)}_{\rho \gamma_1 \gamma_2 \rho} + F^{(2)}_{\rho \gamma_2 \gamma_1 \rho} - F^{(2)}_{\gamma_1 \gamma_2 \rho \rho} - F^{(2)}_{\rho \rho \gamma_1 \gamma_2}\bigr) \nonumber\\
&\quad{}+ \frac{2}{9(d-2)} \delta_{\gamma_1 \gamma_2} \bigl(F^{(2)}_{\rho \alpha \beta \rho} + F^{(2)}_{\rho \beta \alpha \rho} - F^{(2)}_{\alpha \beta \rho \rho} - F^{(2)}_{\rho \rho \alpha \beta}\bigr) \nonumber\\
&\quad{}+ \frac{1}{9(d-2)} \delta_{\gamma_1 \alpha} \bigl(F^{(2)}_{\gamma_2 \beta \rho \rho} + F^{(2)}_{\rho \rho \gamma_2 \beta} - F^{(2)}_{\rho \gamma_2 \beta \rho} - F^{(2)}_{\rho \beta \gamma_2 \rho}\bigr) \nonumber\\
&\quad{}+ \frac{1}{9(d-2)} \delta_{\gamma_2 \alpha} \bigl(F^{(2)}_{\gamma_1 \beta \rho \rho} + F^{(2)}_{\rho \rho \gamma_1 \beta} - F^{(2)}_{\rho \gamma_1 \beta \rho} - F^{(2)}_{\rho \beta \gamma_1 \rho}\bigr) \nonumber\\
&\quad{}+ \frac{1}{9(d-2)} \delta_{\gamma_1 \beta} \bigl(F^{(2)}_{\gamma_2 \alpha \rho \rho} + F^{(2)}_{\rho \rho \gamma_2 \alpha} - F^{(2)}_{\rho \gamma_2 \alpha \rho} - F^{(2)}_{\rho \alpha \gamma_2 \rho}\bigr) \nonumber\\
&\quad{}+ \frac{1}{9(d-2)} \delta_{\gamma_2 \beta} \bigl(F^{(2)}_{\gamma_1 \alpha \rho \rho} + F^{(2)}_{\rho \rho \gamma_1 \alpha} - F^{(2)}_{\rho \gamma_1 \alpha \rho} - F^{(2)}_{\rho \alpha \gamma_1 \rho}\bigr) . \label{conf-obs}
\end{align}

We may also compute the same projection using the fact that $\bG_{\alpha \beta} = N_{\alpha} \cdot N_{\beta}$. This computation is somewhat involved, but results in
\begin{align}
P_{\scalebox{0.2}{\ydiagram{2,2}}}\, N^A_{\gamma_1} N^B_{\gamma_2} \hd_A \hd_B \bG_{\alpha \beta} &{}\= -2\beta_{a \gamma_1 (\alpha} \beta^a_{\beta) \gamma_2} - \frac{2}{3} W_{\gamma_1 (\alpha \beta) \gamma_2} \nonumber\\
&\quad{}+ \frac{2}{3(d-2)} \bigl[
\delta_{\gamma_1 \gamma_2} \bigl(- \IIo^2_{\alpha \beta} + 2 \beta_{a \alpha \rho} \beta^a_{\rho \beta}\bigr)\! + \delta_{\alpha \beta} \bigl(- \IIo^2_{\gamma_1 \gamma_2} + 2 \beta_{a \gamma_1 \rho} \beta^a_{\rho \gamma_2}\bigr) \nonumber\\
&\quad{}
+ \delta_{\gamma_1 (\alpha} \bigl(\IIo^2_{\beta) \gamma_2 } - 2 \beta^a_{\beta) \rho } \beta_{a \rho \gamma_2} \bigr)
+ \delta_{\gamma_2 (\alpha} \bigl(\IIo^2_{\beta) \gamma_1 } - 2\beta^a_{ \beta) \rho} \beta_{a \rho \gamma_1} \bigr) \bigr] . \!\!\!\label{conf-obs-res}
\end{align}

Now, as traces are combined in non-trivial ways with the above expressions, it is fruitful to compute the irreducible components of $F^{(2)}|_{\Lambda}$ independently. To do so, first observe that
\[Tr_{\scalebox{0.2}{\ydiagram{2,2}}}^2 F^{(2)} := \delta^{\alpha \beta} \delta^{\gamma \delta} P_{\scalebox{0.2}{\ydiagram{2,2}}}\, F^{(2)}_{\alpha \beta \gamma \delta} = \frac{2}{3}\bigl(F^{(2)}_{\alpha \alpha \beta \beta} - F^{(2)}_{\alpha \beta \beta \alpha}\bigr) .\]
Evaluating this trace combination on equation~(\ref{conf-obs}), we have that
\[Tr_{\scalebox{0.2}{\ydiagram{2,2}}}^2 N^A_{\gamma_1} N^B_{\gamma_2} \hd_A \hd_B \bG_{\alpha \beta} = \frac{4(3d-2k-4)}{9(d-2)}\bigl(F^{(2)}_{\alpha \alpha \beta \beta} - F^{(2)}_{\alpha \beta \beta \alpha}\bigr) = \frac{2(3d-2k-4)}{3(d-2)} Tr_{\scalebox{0.2}{\ydiagram{2,2}}}^2 F^{(2)} .\]
But we can also take the double trace of equation~(\ref{conf-obs-res}). Doing so yields
\[Tr_{\scalebox{0.2}{\ydiagram{2,2}}}^2\, N^A_{\gamma_1} N^B_{\gamma_2} \hd_A \hd_B \bG_{\alpha \beta} \= -\frac{4(k-1)}{3(d-2)} \IIo^2_{\alpha \alpha} + \frac{2(3d-4k-2)}{3(d-2)} \beta_{a \alpha \beta} \beta^a_{\alpha \beta} - \frac{2}{3} W_{\alpha \beta \beta \alpha} .\]
Thus, we have that
\begin{align*}
Tr_{\scalebox{0.2}{\ydiagram{2,2}}}^2 F^{(2)} \= - \frac{2(k-1)}{3d-2k-4} \IIo^2_{\alpha \alpha} + \frac{3d-4k-2}{3d-2k-4} \beta_{a \alpha \beta} \beta^a_{\alpha \beta} - \frac{d-2}{3d-2k-4} W_{\alpha \beta \beta \alpha} .
\end{align*}

In general, this obstruction \smash{$\mathrm{Tr}_{\scalebox{0.2}{\ydiagram{2,2}}}^2 F^{(2)}$} depends in a non-trivial way on the parallelization of $N \Lambda$ through the appearance of $\beta$, however for special choices of $(d,k)$, we find that this obstruction is a true submanifold invariant that does not depend on the choice of orthonormal frame. This follows because $\IIo_{\alpha}$ and $n_{\alpha}$ both occupy tensor representations of ${\rm O}(k)$, unlike $\beta$. This occurs when $3d-4k-2 = 0$, i.e., for any integer $n \in \mathbb{N}$, then $\mathrm{Tr}_{\scalebox{0.2}{\ydiagram{2,2}}}^2 F^{(2)}$ is a true invariant so long as $(d,k) = (4n+2,3n+1)$.

We next consider the tracefree part of the trace of $F^{(2)}$. As above, it is useful to compute which component this is in terms of $F^{(2)}$ prior to projection. There we have that
\[\bigl[\bigl( \operatorname{t.f.} \circ \mathrm{Tr}_{\scalebox{0.2}{\ydiagram{2,2}}} \bigr) F^{(2)}\bigr]_{\alpha \beta} := \frac{1}{3} F^{(2)}_{(\alpha \beta)_\circ \rho \rho} + \frac{1}{3} F^{(2)}_{\rho \rho (\alpha \beta)_\circ} - \frac{2}{3} F^{(2)}_{\rho (\alpha \beta)_\circ \rho} ,\]
where $\operatorname{t.f.}$ is a projection operator to the trace-free part of a tensor.
Taking this trace of equations~(\ref{conf-obs}) and~(\ref{conf-obs-res}), we find that
\begin{gather*}
\bigl[\bigl( \operatorname{t.f.} \circ \mathrm{Tr}_{\scalebox{0.2}{\ydiagram{2,2}}} \bigr) F^{(2)}\bigr]_{\alpha \beta} \= - \frac{k-2}{3(3d-k-4)} \IIo^2_{(\alpha \beta)_\circ} - \frac{3d-2k-2}{3(3d-k-4)} \beta_{a \rho (\alpha} \beta^a_{\beta)_\circ \rho} \\ \hphantom{\bigl[\bigl( \operatorname{t.f.} \circ Tr_{\scalebox{0.2}{\ydiagram{2,2}}} \bigr) F^{(2)}\bigr]_{\alpha \beta} \=}{}
 - \frac{d-2}{3(3d-k-4)} W_{\rho (\alpha \beta)_{\circ} \rho} .
\end{gather*}
Unlike the double trace obstruction, there are no pairs $(d,k)$ such that this is a true extrinsic submanifold invariant. The final component of $F^{(2)}$, the trace-free part, can be read off from equations~(\ref{conf-obs}) and~(\ref{conf-obs-res}).

Now, given that $F^{(2)}$ occupies the window tensor representation, we have that $F^{(2)}_{\alpha \beta \gamma \delta} = F^{(2)}_{\gamma \delta \alpha \beta}$ and that \smash{$F^{(2)}_{(\alpha \beta \gamma) \delta} = 0$}. It follows that \smash{$F^{(2)}_{\alpha \beta \gamma \delta} + 2 F^{(2)}_{\gamma (\alpha \beta) \delta} = 0$}, and so
\[Tr_{\scalebox{0.2}{\ydiagram{2,2}}}^2 F^{(2)} = \frac{1}{3}F^{(2)}_{\alpha \alpha \beta \beta} \qquad \text{and} \qquad \bigl[\bigl( \operatorname{t.f.} \circ \mathrm{Tr}_{\scalebox{0.2}{\ydiagram{2,2}}} \bigr) F^{(2)}\bigr]_{\alpha \beta} = \frac{2}{3} F^{(2)}_{\rho \rho (\alpha \beta)_{\circ}} .\]
We now summarize all of these results in the following theorem.
\begin{Theorem} \label{ho-conformal-construction}
Let $\Lambda^{d-k} \hookrightarrow \bigl(M^d,\cc\bigr)$ be a parallelized submanifold embedding with $k \neq d-2$. Then there exists a family of defining densities $\sigma_{\alpha}$ for $\Lambda$ that agree modulo terms of order $\sigma^4$ such that
\[\bG_{\alpha \beta} = \delta_{\alpha \beta} + F^{(2)}_{\alpha \beta \gamma_1 \gamma_2} \sigma_{\gamma_1} \sigma_{\gamma_2} ,\]
where $F^{(2)} \in C^\infty M \times \raisebox{3pt}{\scalebox{0.2}{\ydiagram{2,2}}}$, and
\[\operatorname{t.f.} F^{(2)}_{\alpha \beta \gamma_1 \gamma_2} \= \operatorname{t.f.} \biggl(- \beta_{a \gamma_1 (\alpha} \beta^a_{\beta) \gamma_2} - \frac{1}{3} W_{\gamma_1 (\alpha \beta) \gamma_2} \biggr) ,\]
\[F^{(2)}_{\rho \rho (\alpha \beta)_\circ} \= - \frac{k-2}{2(3d-k-4)} \IIo^2_{(\alpha \beta)_\circ} - \frac{3d-2k-2}{2(3d-k-4)} \beta_{a \rho (\alpha} \beta^a_{\beta)_\circ \rho} - \frac{d-2}{2(3d-k-4)} W_{\rho (\alpha \beta)_{\circ} \rho} ,\]
and
\[F^{(2)}_{\alpha \alpha \beta \beta}\= - \frac{6(k-1)}{3d-2k-4} \IIo^2_{\alpha \alpha} + \frac{3(3d-4k-2)}{3d-2k-4} \beta_{a \alpha \beta} \beta^a_{\alpha \beta} - \frac{3(d-2)}{3d-2k-4} W_{\alpha \beta \beta \alpha} .\]
\end{Theorem}
\noindent
We call these defining densities \textit{associated defining densities} for $\Lambda^{d-k} \hookrightarrow \bigl(M^d,g\bigr)$.

